% Options for packages loaded elsewhere
\PassOptionsToPackage{unicode}{hyperref}
\PassOptionsToPackage{hyphens}{url}
\PassOptionsToPackage{dvipsnames,svgnames,x11names}{xcolor}
\documentclass[
  letterpaper,
  11pt]{article}
\usepackage{xcolor}
\usepackage[margin=1in]{geometry}
\usepackage{amsmath,amssymb}
\setcounter{secnumdepth}{-\maxdimen} % remove section numbering
\usepackage{iftex}
\ifPDFTeX
  \usepackage[T1]{fontenc}
  \usepackage[utf8]{inputenc}
  \usepackage{textcomp} % provide euro and other symbols
\else % if luatex or xetex
  \usepackage{unicode-math} % this also loads fontspec
  \defaultfontfeatures{Scale=MatchLowercase}
  \defaultfontfeatures[\rmfamily]{Ligatures=TeX,Scale=1}
\fi
\usepackage{lmodern}
\ifPDFTeX\else
  % xetex/luatex font selection
  \setmainfont[]{Latin Modern Roman}
  \setmonofont[]{Latin Modern Mono}
  \setmathfont[]{Latin Modern Math}
\fi
% Use upquote if available, for straight quotes in verbatim environments
\IfFileExists{upquote.sty}{\usepackage{upquote}}{}
\IfFileExists{microtype.sty}{% use microtype if available
  \usepackage[]{microtype}
  \UseMicrotypeSet[protrusion]{basicmath} % disable protrusion for tt fonts
}{}
\makeatletter
\@ifundefined{KOMAClassName}{% if non-KOMA class
  \IfFileExists{parskip.sty}{%
    \usepackage{parskip}
  }{% else
    \setlength{\parindent}{0pt}
    \setlength{\parskip}{6pt plus 2pt minus 1pt}}
}{% if KOMA class
  \KOMAoptions{parskip=half}}
\makeatother
\usepackage{longtable,booktabs,array}
\newcounter{none} % for unnumbered tables
\usepackage{calc} % for calculating minipage widths
% Correct order of tables after \paragraph or \subparagraph
\usepackage{etoolbox}
\makeatletter
\patchcmd\longtable{\par}{\if@noskipsec\mbox{}\fi\par}{}{}
\makeatother
% Allow footnotes in longtable head/foot
\IfFileExists{footnotehyper.sty}{\usepackage{footnotehyper}}{\usepackage{footnote}}
\makesavenoteenv{longtable}
\setlength{\emergencystretch}{3em} % prevent overfull lines
\providecommand{\tightlist}{%
  \setlength{\itemsep}{0pt}\setlength{\parskip}{0pt}}
% =============================================================================
% Custom preamble for simplicial_vector_calculus (pandoc -> XeLaTeX)
% Pulled in via `--include-in-header` by build.ps1 / build.sh.
%
% This file is injected into pandoc's default LaTeX template at the
% `header-includes` slot, which runs BEFORE pandoc's own \usepackage{hyperref}
% and \usepackage{unicode-math}. We therefore:
%   * load `mathtools` here (so it precedes unicode-math, per its manual),
%   * defer any \hypersetup{...} with \AtBeginDocument{...}.
% We only ADD or tune; we do not reload packages pandoc already loads
% (amsmath, amssymb, amsfonts, hyperref, xcolor, fontspec, unicode-math).
% =============================================================================

% ---- Math / symbols (mathtools must precede unicode-math) ------------------
\usepackage{mathtools}
\usepackage{amsthm}

% ---- Fix \overbrace/\underbrace under XeLaTeX + unicode-math ----------------
% unicode-math breaks amsmath's \over/\underbrace: their \up/\downbracefill use
% cmex math-chars that remap to stray glyphs (z, {, }) padded by \hrule leaders
% (black bars). Route them to mathtools' rule-based \over/\underbracket, which
% render correctly and accept the identical \underbrace{...}_{label} /
% \overbrace{...}^{label} syntax. Square corners instead of curly (cosmetic).
% Deferred to \AtBeginDocument so it runs after pandoc loads unicode-math.
% unicode-math overwrites \over/\underbracket with a no-optional-argument form,
% so pass no [thickness][height] optionals (they would leak as literal text).
\AtBeginDocument{%
  \renewcommand{\underbrace}[1]{\underbracket{#1}}%
  \renewcommand{\overbrace}[1]{\overbracket{#1}}%
}

% ---- Typography -------------------------------------------------------------
\usepackage{microtype}
\usepackage{booktabs}

% ---- Title-page vertical compression ---------------------------------------
% Pull the title up so the abstract + sharing box fit on page 1.
\usepackage{titling}
\setlength{\droptitle}{-3em}

% ---- Force the table of contents onto its own page -------------------------
% Decouples page 1 layout from word-count fluctuations: title + abstract
% stay on page 1, Contents starts at the top of page 2.
\let\origtableofcontents\tableofcontents
\renewcommand{\tableofcontents}{\clearpage\origtableofcontents}

% ---- Figures ----------------------------------------------------------------
\usepackage{tikz}
\usepackage{float}   % enables the [H] float specifier (lock figure in place)

% ---- Callout / banner boxes ------------------------------------------------
\usepackage{mdframed}

% ---- Hebrew letter support (for the trilogy labels ל1, ל2, ל3) -----------
% Latin Modern Roman lacks Hebrew glyphs, so xelatex would otherwise emit
% "missing character" warnings and leave the slots blank. We define a
% \heb{...} command that switches to a system font with Hebrew coverage
% (Times New Roman, which ships with Windows) just for the wrapped text.
%
% \heb is declared *robust* at preamble time so packages that
% protected-edef the title argument (titling, hyperref) do not try to
% expand \heb when its real definition is not yet in place. \hebrewfont
% is a no-op stub at preamble time; the real fontspec-based font family
% is installed inside \AtBeginDocument{}, since fontspec is loaded by
% pandoc only after the header-includes.
%
% Accessibility: the glyph ל (Hebrew lamed) confuses PDF-to-text pipelines
% (pdftotext, Zenodo's OCR, screen readers, search indexers), which often
% drop it or replace it with a placeholder. We wrap the glyph in an
% accsupp ActualText annotation so the *extracted* text reads "L" while
% the *rendered* glyph remains ל. This makes "ל5" extract as "L5" — the
% intended Latin transliteration of the labels — without changing the
% printed paper at all.
\usepackage{accsupp}
\providecommand{\hebrewfont}{}
% Optional first argument is the ActualText used by PDF text extraction; it
% defaults to "L" for the ל (lamed / interior) series. The פ (pe / exterior)
% series passes [P], so \heb[P]{פ} renders the glyph פ but extracts as "P".
% The default keeps every existing \heb{ל} call extracting as "L" unchanged.
\DeclareRobustCommand{\heb}[2][L]{%
  \ifmmode
    \text{\BeginAccSupp{ActualText=#1}{\hebrewfont #2}\EndAccSupp{}}%
  \else
    \BeginAccSupp{ActualText=#1}{\hebrewfont #2}\EndAccSupp{}%
  \fi
}
\AtBeginDocument{%
  \newfontfamily\HebrewSysFont{Times New Roman}[Scale=MatchLowercase]%
  \renewcommand{\hebrewfont}{\HebrewSysFont}%
}
% Let LaTeX stretch lines a bit further when it can't find a clean break;
% keeps overfull hboxes out of display math and URL-heavy lines.
\setlength{\emergencystretch}{3em}
% Make long URLs breakable at hyphens (applied before hyperref loads).
\PassOptionsToPackage{hyphens}{url}

% ---- Section heading polish -------------------------------------------------
\usepackage{titlesec}
\titleformat{\section}
  {\normalfont\Large\bfseries}{\thesection}{0.75em}{}
\titleformat{\subsection}
  {\normalfont\large\bfseries}{\thesubsection}{0.6em}{}
\titleformat{\subsubsection}
  {\normalfont\normalsize\bfseries}{\thesubsubsection}{0.5em}{}

% ---- Hyperref polish (hyperref is loaded later by pandoc; defer) ------------
% PDF metadata (pdftitle / pdfauthor / pdfsubject / pdfkeywords) is set by
% pandoc's default LaTeX template from the YAML front matter (title, author,
% subject, keywords) — we intentionally do not override it here so the same
% preamble can serve any paper in the repo.

% ---- Theorem-like environments (available if source ever migrates) ---------
% The markdown source uses `**Theorem 4.1 (title).**` bold-header paragraphs
% rather than native environments, so these are currently unused by the body
% but are kept here so a future migration does not require a preamble change.
\theoremstyle{plain}
\newtheorem{theorem}{Theorem}[section]
\newtheorem{proposition}[theorem]{Proposition}
\newtheorem{lemma}[theorem]{Lemma}
\newtheorem{corollary}[theorem]{Corollary}
\theoremstyle{definition}
\newtheorem{definition}[theorem]{Definition}
\theoremstyle{remark}
\newtheorem{remark}[theorem]{Remark}
\newtheorem{example}[theorem]{Example}
\providecommand{\R}{\mathrm{R}}
\usepackage{bookmark}
\IfFileExists{xurl.sty}{\usepackage{xurl}}{} % add URL line breaks if available
\urlstyle{same}
\hypersetup{
  pdftitle={Exact-Arithmetic Verification of the Spectral Realization by Zeta-Cycles},
  pdfauthor={Leonardo Murillo Montero (Independent researcher --- \href{mailto:lmurillo@avionyx.com}{\nolinkurl{lmurillo@avionyx.com}})},
  colorlinks=true,
  linkcolor={Maroon},
  filecolor={Maroon},
  citecolor={Blue},
  urlcolor={Blue},
  pdfcreator={LaTeX via pandoc}}

\title{Exact-Arithmetic Verification of the Spectral Realization by Zeta-Cycles: the Simple-Even Condition, the Prolate Proximity, and the Eigenvalue Ladder of \(QW_\lambda\)}
\author{Leonardo Murillo Montero\newline\textit{Independent researcher — lmurillo@avionyx.com}}
\date{June 2026}

\begin{document}
\maketitle
\begin{abstract}
This note reports an independent, high-precision implementation of the semilocal Weil quadratic form \(QW_\lambda\) of Connes--Consani--Moscovici (arXiv:2511.22755, hereafter {[}CCM{]}), together with exact-arithmetic evidence bearing directly on the two missing steps identified in {[}CCM{]}, Section 8. The contributions are: \textbf{(1)} a reduction of the matrix of \(QW_\lambda^N\) in the basis \(V_n\) to \(O(N)\) one-dimensional integrals of elementary functions, permitting fast evaluation at 40-digit precision; \textbf{(2)} verification of the \emph{simple-even condition} ({[}CCM{]}, first missing step) at the windows \(\lambda = 1.3,\ 1.5,\ 2.0\), with both parity sectors computed so that evenness is verified rather than imposed; \textbf{(3)} a direct quantitative measurement of the proximity of the prolate ansatz \(k_\lambda\) to the true ground vector \(\xi_\lambda\) ({[}CCM{]}, second missing step): the normalized Rayleigh excess is \(r/\epsilon_\lambda \approx 0.14\), and the resulting min--max bound \(\|\tilde k_\lambda - \xi_\lambda\| \le \sqrt{2r/(\mu_1-\epsilon)}\) is verified on the data; \textbf{(4)} a quantification, in the exact spectra, of the sharply graded \emph{eigenvalue ladder} whose existence is indicated in {[}CCM{]}, Section 8 (indication (3)): the ratio \(\mu_1/\epsilon_\lambda\) of the first excited even eigenvalue to the ground eigenvalue grows from \(\approx 37\) to \(\approx 1.8\times10^3\) to \(\approx 2\times10^5\) across the three windows, with the first excited even eigenvector lying in the span of the \(\mathcal E\)-images of the prolate cells \(n = 0,4,8\) with squared-projection deficit \(\approx 3\times10^{-6}\) --- so that the spectral gap relevant to the convergence strategy of {[}CCM{]}, Section 7, is exponentially larger than the ground eigenvalue; \textbf{(5)} an independent reproduction of the identity \(\epsilon_\lambda \approx 1-\chi(\lambda)\) ({[}CCM{]}, Figure 4) over thirteen orders of magnitude, and of the spectral realization itself: with the primes \(2, 3\) (and the prime power \(4\)) only, the zeros of \(\widehat\xi\) computed from our matrix are real to the working-precision floor (\(10^{-44}\) at 40 digits, falling to \(10^{-66}\) at 60 digits) and reproduce \(\gamma_1\) to \(2.3\times10^{-9}\). All results are packaged as five deterministic Python certificates (attached, and mirrored at the repository linked below), including a stability audit (truncation refinement, quadrature doubling, precision doubling, root-finder basin), which run in minutes on a laptop.
\end{abstract}

\emph{Note, version 2 --- prepared for the authors of arXiv:2511.22755}

\subsection{1. Introduction: what this note contributes}\label{introduction-what-this-note-contributes}

The paper {[}CCM{]} constructs spectral triples from rank-one perturbations \(\widetilde D\) of the scaling operator on \([\lambda^{-1},\lambda]\), governed by the restriction \(QW_\lambda^N\) of the Weil quadratic form to the span \(E_N\) of the \(2N+1\) lowest modes \(V_n\). Its main theorem produces, from the lowest eigenvector \(\xi\) of \(QW_\lambda^N\) --- \emph{assumed simple and even, normalized by \(\delta_N(\xi)=1\)} --- a regularized determinant \(\det\nolimits_\infty(\widetilde D - z) = -i\lambda^{-iz}\widehat\xi(z)\) whose zeros are all real ({[}CCM{]}, Theorem 1.1, building on the engine of Connes--van Suijlekom {[}CS{]}), and the numerical sections of {[}CCM{]} and of the predecessor paper {[}CC-cycles{]} show these zeros converging to the nontrivial zeros of \(\zeta\) with striking accuracy. Section 8 of {[}CCM{]} states the two steps still missing from this strategy, in the authors' words:

\begin{quote}
``The first is that, in order to apply Theorem 5.10 to the Weil quadratic form \(QW_\lambda\), one must prove that its smallest eigenvalue---whose existence is ensured by Theorem 3.6---is simple and that its corresponding eigenvector \(\xi_\lambda\) is even. The second step is to establish that \(k_\lambda\) provides a sufficiently accurate approximation to (a scalar multiple of) \(\xi_\lambda\), in order to justify the convergence of the zeros of \(\widehat\xi_\lambda\) towards the non-trivial zeros of \(\zeta(\frac12+is)\).''
\end{quote}

This note contributes computational theorems-of-record on both, obtained from an implementation that is independent of the authors' (different reduction of the matrix elements, different prolate solver, arbitrary-precision arithmetic throughout), so that agreement constitutes genuine verification. We are aware of two other independent numerical engagements with this circle of results, which the present note complements rather than repeats: Groskin (arXiv:2605.20224) has implemented the Connes--van Suijlekom Galerkin matrix at cutoffs \(c = \lambda^2 \ge 13\) to very high precision, with the even sector imposed by the choice of basis; and Śliwiński (arXiv:2601.12133) has computed low-precision spectra of the operators \(D^{(\lambda,N)}_{\log}\) against the first thousand zeros. Neither covers the semilocal windows \(\lambda \le 2\), computes the odd sector (so that parity is \emph{verified}, not assumed), measures the prolate-ansatz proximity, or resolves the eigenvalue ladder --- those are the regimes and observables reported here. Beyond verification, the exact spectra quantify structure that {[}CCM{]}, Section 8 (indication (3)) describes qualitatively --- the eigenvalue ladder of Section 4 --- and this, we believe, materially improves the prospects of step (ii): the gap \(\mu_1 - \epsilon_\lambda\) entering any perturbative comparison of \(k_\lambda\) with \(\xi_\lambda\) is measured to be not merely nonvanishing but \emph{exponentially large relative to} \(\epsilon_\lambda\).

We follow the notation of {[}CCM{]} throughout: \(L = 2\log\lambda\); \(U_n(x) = L^{-1/2}e^{2\pi inx/L}\) on \([0,L]\); \(V_n = \kappa(U_n)\); \(\mathcal E(f)(u) = u^{1/2}\sum_{k\ge1}f(ku)\); \(h_{n,\lambda}\) the prolate eigenfunctions of \(PW_\lambda = -\partial_x((\lambda^2-x^2)\partial_x) + (2\pi\lambda x)^2\); \(\chi(\lambda)\) the eigenvalue of the Fourier compression; \(k_\lambda = \mathcal E(h_\lambda)\) with \(h_\lambda\) the vanishing-integral combination of \(h_{0,\lambda}, h_{4,\lambda}\).

\subsection{2. The exact matrix: a per-frequency reduction}\label{the-exact-matrix-a-per-frequency-reduction}

Our starting point is the representation of \(QW(V_n, V_m)\) through the distribution \(\Psi^\sharp\) on \([1,\infty)\) ({[}CCM{]}, Lemma 3.1 and Proposition 3.2): \(QW(V_n,V_m) = \Psi^\sharp(F_{nm})\) with \(F_{nm}(x) = q(U_n,U_m)(\log x)\), where ({[}CCM{]}, Lemma 2.3)
\[
q(U_n,U_m)(y) = \frac{\sin(\omega_m y) - \sin(\omega_n y)}{\pi(n-m)}\ \ (n\neq m), \qquad
q(U_n,U_n)(y) = 2\Big(1 - \frac{y}{L}\Big)\cos(\omega_n y), \qquad \omega_k := \frac{2\pi k}{L}.
\]

The observation that makes exact evaluation cheap is that \(\Psi^\sharp\) is \emph{linear} in \(q\) and the off-diagonal \(q\)'s are differences of single-frequency sines. Hence the entire \((2N+1)^2\) matrix is generated by \(2N+1\) scalars: setting, in the variable \(y = \log x\),
\[
S_k := \underbrace{\int_0^L \sin(\omega_k y)\,2\cosh(y/2)\,dy}_{A_k\ (\text{from }W_{0,2}^\sharp)}
\;-\;\underbrace{\int_0^L \frac{e^{y/2}\sin(\omega_k y)}{2\sinh y}\,dy}_{B_k\ (\text{from }W_\R^\sharp)}
\;-\;\underbrace{\sum_{p^j\le e^L}\frac{\log p}{p^{j/2}}\,\sin(\omega_k\, j\log p)}_{P_k\ (\text{from }W_p^\sharp)},
\]
one has
\[
QW(V_n,V_m) \;=\; \frac{S_m - S_n}{\pi(n-m)} \qquad (n\neq m),
\]
with an analogous expression for the diagonal from \(q(U_n,U_n)\); the \(F(1)\)-term of \(W_\R^\sharp\) contributes, beyond the window, the closed form \(\tfrac12\log\tanh(L/2)\). (The \(B_k\) integrand extends continuously by \(\omega_k/2\) at \(y=0\).) All integrands are smooth and elementary; we evaluate them by adaptive Gauss--Legendre quadrature in \texttt{mpmath} at \(40\) significant digits. The cost of the full matrix at \(N = 12\) is a few seconds.

\textbf{Validation.} The diagonal entries \(QW(V_0,V_0)\) and \(QW(V_1,V_1)\) at \(\lambda = 1.5\) were checked against the zero side of the explicit formula, \(\sum_\rho\) over the first \(200\) zeros with the standard tail estimate: agreement within the truncation allowance. Doubling the quadrature subdivision (\(16 \to 32\) Gauss--Legendre panels) leaves every reported eigenvalue unchanged to ten significant digits, and raising the working precision from \(40\) to \(60\) digits changes them only at the \(10^{-29}\) level --- the quoted numbers are quadrature- and precision-exact. Independently, the spectral-realization computation of Section 6 --- which exercises every entry of the matrix through the eigenvector --- reproduces \(\gamma_1\) to \(2.3\times10^{-9}\), a far more stringent end-to-end check. (Certificates: \texttt{probe\_ccm\_exact\_matrix.py}, CHECK A; \texttt{attack\_v2\_hardening.py}, ATTACK 1.)

\subsection{3. The simple-even condition (first missing step)}\label{the-simple-even-condition-first-missing-step}

For each window we computed the full spectrum of \(QW^N_\lambda\) on \(E_N\) and separated the even sector (\(e_0 = V_0\), \(e_k = (V_k + V_{-k})/\sqrt2\)) from the odd. The low spectra are:

{\def\LTcaptype{none} % do not increment counter
\begin{longtable}[]{@{}
  >{\raggedright\arraybackslash}p{(\linewidth - 8\tabcolsep) * \real{0.2000}}
  >{\raggedright\arraybackslash}p{(\linewidth - 8\tabcolsep) * \real{0.2000}}
  >{\raggedright\arraybackslash}p{(\linewidth - 8\tabcolsep) * \real{0.2000}}
  >{\raggedright\arraybackslash}p{(\linewidth - 8\tabcolsep) * \real{0.2000}}
  >{\raggedright\arraybackslash}p{(\linewidth - 8\tabcolsep) * \real{0.2000}}@{}}
\toprule\noalign{}
\begin{minipage}[b]{\linewidth}\raggedright
\(\lambda\)
\end{minipage} & \begin{minipage}[b]{\linewidth}\raggedright
\(N\)
\end{minipage} & \begin{minipage}[b]{\linewidth}\raggedright
prime powers in window
\end{minipage} & \begin{minipage}[b]{\linewidth}\raggedright
even spectrum (lowest four)
\end{minipage} & \begin{minipage}[b]{\linewidth}\raggedright
odd spectrum (lowest two)
\end{minipage} \\
\midrule\noalign{}
\endhead
\bottomrule\noalign{}
\endlastfoot
\(1.3\) & \(10\) & none & \(2.46\times10^{-2}\), \(0.905\), \(1.47\), \(1.83\) & \(0.375\), \(1.22\) \\
\(1.5\) & \(10\) & \(2\) & \(1.567\times10^{-4}\), \(0.280\), \(0.904\), \(1.36\) & \(1.44\times10^{-2}\), \(0.898\) \\
\(2.0\) & \(12\) & \(2,3,4\) & \(1.375\times10^{-12}\), \(2.753\times10^{-7}\), \(3.76\times10^{-3}\), \(0.568\) & \(8.75\times10^{-10}\), \(3.80\times10^{-5}\) \\
\end{longtable}
}

\textbf{In each case the lowest eigenvalue of \(QW^N_\lambda\) is simple, its eigenvector is even, and the odd minimum lies strictly above it.} At \(\lambda = 2.0\), \(N = 12\), the separation between the even ground state and the odd minimum is a factor \(6.4\times10^2\); between the even ground state and the next \emph{even} eigenvalue, a factor \(2.0\times10^5\).

A truncation disclosure is in order, since these are spectra of the finite sections \(QW^N_\lambda\) --- which are precisely the objects entering {[}CCM{]}, Theorem 1.1. The low eigenvalues decrease monotonically under refinement (consistent with Cauchy interlacing of nested principal submatrices) and the smallest are not yet \(N\)-converged in their leading digits: at \(\lambda = 2.0\), refining \(N = 12 \to 14 \to 16\) moves the ground eigenvalue \(1.375\times10^{-12} \to 1.108\times10^{-12} \to 1.056\times10^{-12}\), the first excited even eigenvalue \(2.753\times10^{-7} \to 2.717\times10^{-7} \to 2.464\times10^{-7}\), and the odd minimum \(8.75\times10^{-10} \to 7.56\times10^{-10} \to 6.47\times10^{-10}\); at \(\lambda = 1.5\) the ground eigenvalue moves by \(-2.6\%\) (\(N{+}2\)) and \(-4.5\%\) (\(N{+}4\)); at \(\lambda = 1.3\) by \(-0.3\%\). The \emph{qualitative} verdict --- lowest eigenvalue simple, eigenvector even, odd minimum strictly above, ladder ratios stable in order of magnitude (\(\mu_1/\epsilon = 2.0\)--\(2.5\times10^5\) at \(\lambda = 2\) across \(N = 12,14,16\)) --- holds at every truncation tested. The hypothesis of {[}CCM{]}, Theorem 1.1 is thus verified by exact arithmetic at three windows spanning the prime-free and semilocal regimes, at every computed truncation. (Certificates: \texttt{probe\_ccm\_exact\_matrix.py}, CHECK B; \texttt{attack\_v2\_hardening.py}, ATTACK 1.)

\subsection{\texorpdfstring{4. The eigenvalue ladder, and the identity \(\epsilon_\lambda \approx 1 - \chi(\lambda)\)}{4. The eigenvalue ladder, and the identity \textbackslash epsilon\_\textbackslash lambda \textbackslash approx 1 - \textbackslash chi(\textbackslash lambda)}}\label{the-eigenvalue-ladder-and-the-identity-epsilon_lambda-approx-1---chilambda}

Two features of the table above seemed to us worth isolating and quantifying, because they bear directly on the feasibility of the second missing step. (That the low spectrum has such structure is indicated qualitatively in {[}CCM{]}, Section 8, indication (3); the measurements below are, to our knowledge, the first.)

\textbf{(a) The gap is exponentially large relative to the ground eigenvalue.} The ratio \(\mu_1/\epsilon_\lambda\) (first excited even / ground) grows as
\[
\lambda = 1.3\ (N{=}10):\ 37; \qquad \lambda = 1.5\ (N{=}10):\ 1.79\times10^{3}; \qquad \lambda = 2.0\ (N{=}12):\ 2.00\times10^{5},
\]
the last stable in order of magnitude (\(2.3\)--\(2.5\times10^5\)) under \(N \to 14, 16\). Moreover the structure of the excited state is exactly what the prolate picture predicts: the first excited even eigenvector of the exact matrix lies in \(\mathrm{span}\{\mathcal E h_{0,\lambda}, \mathcal E h_{4,\lambda}, \mathcal E h_{8,\lambda}\}\) with projection \(0.999997\) --- squared-projection deficit \(3.2\times10^{-6}\), stable under doubling of both discretization grids and under \(N \to 14\), while twenty random three-dimensional subspaces of the thirteen-dimensional even sector give projections of mean \(0.40\) and maximum \(0.76\) --- and in the span of the first two cells alone with projection only \(0.59\) (\(0.41\) in the full-space/raw-cell convention of the certificate): its \(h_8\)-content is essential. In other words, \emph{the low spectrum of \(QW_\lambda\) is organized as the \(\mathcal E\)-image of the tower of Fourier-invariant prolate cells} \(n \equiv 0 \ (\mathrm{mod}\ 4)\), with the \(k\)-th eigenvalue tracking the \(k\)-th compression defect \(1-\chi_{4k}(\lambda)\), whose Fuchs ratios \((1-\chi_{4(k+1)})/(1-\chi_{4k}) \sim \mathrm{const}\cdot(2\pi\lambda^2)^4\) account for the measured growth. A min--max consequence we record explicitly: compressing the exact \(\lambda=2\), \(N=12\) matrix to the two-dimensional span of the first two cells yields eigenvalues \(1.59\times10^{-12}\) and \(3.93\times10^{-7}\), so by Courant--Fischer \(\mu_1 \le 3.93\times10^{-7}\) --- the gap structure is a theorem at this window, not only a computation (and the measured \(\mu_1 = 2.46\)--\(2.75\times10^{-7}\) at \(N = 12\)--\(16\) is consistent with it). (Certificates: \texttt{probe\_fuchs\_second\_rung.py}; \texttt{attack\_v2\_hardening.py}, ATTACK 2.)

\textbf{(b) \(\epsilon_\lambda\) tracks the Fuchs asymptotics of \(1-\chi(\lambda)\) over thirteen decades.} Using the asymptotic quoted in {[}CCM{]} (from Fuchs, for \(n=4\)),
\[
1-\chi(\lambda)\ \sim\ \tfrac{2^{14}}{3}\sqrt2\,\pi^5\,e^{-4\pi\lambda^2 + 9\log\lambda},
\]
we find \(\epsilon_\lambda / (1-\chi)_{\mathrm{Fuchs}} = 3.3\) at \(\lambda=1.5\) (\(N{=}10\); \(3.1\) at \(N{=}14\)) and \(7.7\) at \(\lambda = 2.0\) (\(N{=}12\); \(6.2\) at \(N{=}14\), \(5.9\) at \(N{=}16\)), while the quantities themselves fall from \(10^{-4}\) to \(10^{-12}\). The truncation is the dominant uncertainty in these ratios; the stable statement is the order-unity tracking across thirteen decades. This reproduces, by an independent route, the phenomenon displayed in {[}CCM{]}, Figure 4, and supports reading \(\epsilon_\lambda\) as the Fourier-compression defect of the ground cell. We also reproduced, from a banded finite-difference diagonalization of \(PW_\lambda\) (independent of any prolate special-function library), the Meixner--Schäfke estimate of {[}CCM{]}, Lemma 7.2: \(\max_{[-\lambda,\lambda]}|h_{n,\lambda}-h_n| = O(\lambda^{-2})\), with fitted rates \(\lambda^{-2.16}\) (\(n=0\)) and \(\lambda^{-2.56}\) (\(n=4\)); and the convergence \(\widehat{k_\lambda}\to\Xi\) of {[}CCM{]}, Lemma 7.3, measured end-to-end at the empirical rate \(\approx\lambda^{-2.3}\) --- faster than the \(O(\lambda^{-1/2-\alpha})\) rate recorded for this convergence in the survey {[}L{]}, Section 6.5. The Poisson symmetry \(k(u)=k(u^{-1})\) holds in our build to \(1.4\times10^{-15}\), and its finite-\(\lambda\) breakage (the discrepancy of \(h_\lambda\) from exact Fourier invariance) measures \(2.5\)--\(3\) orders of magnitude below the \(\lambda^{-2}\) approximation floor --- consistent with the exponential smallness of \(1-\chi(\lambda)\). (Certificate: \texttt{probe\_ccm\_seam\_weld.py}.)

\subsection{\texorpdfstring{5. The proximity of \(k_\lambda\) to \(\xi_\lambda\) (second missing step)}{5. The proximity of k\_\textbackslash lambda to \textbackslash xi\_\textbackslash lambda (second missing step)}}\label{the-proximity-of-k_lambda-to-xi_lambda-second-missing-step}

Let \(\tilde k_\lambda\) denote the \(L^2(d^*u)\)-normalized, even-projected window function built from the prolate ansatz, and let \(r := QW_\lambda(\tilde k_\lambda) - \epsilon_\lambda\) be its Rayleigh excess. At \(\lambda = 1.5\) (\(N = 10\)) the exact matrix gives
\[
QW_\lambda(\tilde k_\lambda) = 1.78\times10^{-4}, \qquad \epsilon_\lambda = 1.567\times10^{-4}, \qquad r = 2.12\times10^{-5}, \qquad r/\epsilon_\lambda = 0.135 .
\]
Since \(\tilde k_\lambda\) is built from finite-difference prolates, we audited the construction error: refining the finite-difference grid (\(2001 \to 4001 \to 8001\) points) and the \(\mathcal E\)-image quadrature grid (\(3001 \to 6001\)) moves \(r/\epsilon_\lambda\) only within \([0.133,\ 0.137]\), the discretized value increasing toward the fine-grid limit \(\approx 0.137\) --- so the construction contributes at the \(2\%\) level to this ratio. We note that each quoted \(QW_\lambda(\tilde k)\) is the exact Rayleigh value of an explicitly constructed vector, hence individually a rigorous upper bound for \(\epsilon_\lambda\); the discretization uncertainty affects only the identification of \(\tilde k\) with the ideal prolate ansatz. By the standard min--max inequality, for any unit vector \(k\) and simple ground pair \((\epsilon, \xi)\) with even-sector gap \(\mu_1\):
\[
\big\|\,k - \langle\xi, k\rangle\xi\,\big\|^2 \ \le\ \frac{QW_\lambda(k) - \epsilon}{\mu_1 - \epsilon}.
\]
On the exact data this reads \(\|\tilde k_\lambda - \xi_\lambda\| = 0.0066 \le \sqrt{2r/(\mu_1-\epsilon)} = 0.0123\): the bound holds with a factor-two margin, and the actual distance is below \(10^{-2}\) already at the smallest semilocal window. Since (Section 4) the gap ratio \(\mu_1/\epsilon_\lambda\) grows at a rate consistent with \((2\pi\lambda^2)^4\) per prolate cell while \(r\) tracks \(\epsilon_\lambda\) itself, the closeness available to the strategy of {[}CCM{]}, Section 7 appears to \emph{improve} exponentially with \(\lambda\): the quantity \(\sqrt{2r/(\mu_1-\epsilon)}\), which controls the eigenvector transfer, is measured at \(1.2\times10^{-2}\) at \(\lambda=1.5\) and is predicted by the ladder to fall like a power of \(e^{-\lambda^2}\) thereafter. For the Mellin side we verified, at four points of the strip including \(\Im z = 0.45\), the elementary transfer bound \(|\widehat{(\tilde k-\xi)}(z)| \le \|\tilde k - \xi\|\cdot W(\beta,\lambda)\) with \(W(\beta,\lambda) = ((\lambda^{2\beta} - \lambda^{-2\beta})/2\beta)^{1/2}\) --- so that eigenvector closeness transfers to the determinant with only polynomially growing constants. (Certificates: \texttt{probe\_quantified\_weld\_e2e.py}, CHECKs A--B; \texttt{attack\_v2\_hardening.py}, ATTACK 4.)

The proximity of \(k_\lambda\) to \(\xi_\lambda\) as such is not new: it is reported as numerical evidence in {[}CC-cycles{]}, Section 3, and {[}CCM{]}, Section 8, indication (3), and the variational use of \(k_\lambda\) originates with the authors. What we believe is new here is the quantitative form --- the Rayleigh excess \(r\), the normalized ratio \(r/\epsilon_\lambda\), and the verified min--max transfer to \(\|\tilde k_\lambda - \xi_\lambda\|\) --- these being the specific quantities that a perturbative proof of step (ii) would need to control, measured for the first time, and uniformly favorable to it.

\subsection{6. The spectral realization from the independent build}\label{the-spectral-realization-from-the-independent-build}

Finally, as an end-to-end exercise of the whole architecture, we computed \(\widehat\xi(z) = \sum_n \xi_n \widehat{V_n}(z)\) from the ground eigenvector of our exact matrix and located its zeros. In accordance with {[}CCM{]}, Theorem 1.1 (iii), the zeros are real: the imaginary parts returned by the root-finder are at the \(10^{-44}\) level at our working precision of 40 digits. That this is genuinely a precision floor --- not a residual imaginary part --- is confirmed by repeating the computation at 60 digits: the imaginary parts fall to the \(10^{-66}\) level while the real parts are unchanged to 46 digits. The root-finder converges to the same root from four distinct starting points per zero (real-part spread \(< 10^{-43}\)), so the values are not artifacts of the search basin. Their positions, against the actual ordinates \(\gamma_k\) of \(\zeta\) (at the truncations of Section 3):

{\def\LTcaptype{none} % do not increment counter
\begin{longtable}[]{@{}
  >{\raggedright\arraybackslash}p{(\linewidth - 4\tabcolsep) * \real{0.3333}}
  >{\raggedright\arraybackslash}p{(\linewidth - 4\tabcolsep) * \real{0.3333}}
  >{\raggedright\arraybackslash}p{(\linewidth - 4\tabcolsep) * \real{0.3333}}@{}}
\toprule\noalign{}
\begin{minipage}[b]{\linewidth}\raggedright
\end{minipage} & \begin{minipage}[b]{\linewidth}\raggedright
\(\lambda = 1.5\) (prime \(2\))
\end{minipage} & \begin{minipage}[b]{\linewidth}\raggedright
\(\lambda = 2.0\) (primes \(2,3\); prime power \(4\))
\end{minipage} \\
\midrule\noalign{}
\endhead
\bottomrule\noalign{}
\endlastfoot
\(|\hat z_1 - \gamma_1|\) & \(4.3\times10^{-2}\) & \(2.3\times10^{-9}\) \\
\(|\hat z_2 - \gamma_2|\) & \(0.73\) & \(6.5\times10^{-7}\) \\
\(|\hat z_3 - \gamma_3|\) & \(3.0\) & \(1.4\times10^{-5}\) \\
\end{longtable}
}

(At \(N = 14\) the \(\lambda = 2.0\) agreement with \(\gamma_1\) improves to \(1.9\times10^{-9}\); the table is at \(N = 12\).) Seven orders of magnitude of improvement on \(\gamma_1\) in half a unit of \(\lambda\) --- the \(e^{-4\pi\lambda^2}\) pace, visible. We confirm, from an implementation sharing no code or method with the authors', the phenomenon that motivates {[}CCM{]}; for the complementary high-cutoff regime \(c = \lambda^2 \ge 13\) see also the independent reproduction of Groskin (arXiv:2605.20224). (Certificates: \texttt{probe\_quantified\_weld\_e2e.py}, CHECK C; \texttt{attack\_v2\_hardening.py}, ATTACK 3.)

\subsection{7. Certificates and reproducibility}\label{certificates-and-reproducibility}

The five attached scripts are deterministic, self-contained, and run on stock Python (\$\ge\$3.11 with \texttt{numpy}, \texttt{scipy}, \texttt{mpmath}); total runtime is under thirty minutes on a laptop. Each prints a PASS/FAIL verdict per check (the stability audit prints measured drifts) and a final result line.

{\def\LTcaptype{none} % do not increment counter
\begin{longtable}[]{@{}
  >{\raggedright\arraybackslash}p{(\linewidth - 4\tabcolsep) * \real{0.3333}}
  >{\raggedright\arraybackslash}p{(\linewidth - 4\tabcolsep) * \real{0.3333}}
  >{\raggedright\arraybackslash}p{(\linewidth - 4\tabcolsep) * \real{0.3333}}@{}}
\toprule\noalign{}
\begin{minipage}[b]{\linewidth}\raggedright
script
\end{minipage} & \begin{minipage}[b]{\linewidth}\raggedright
contents
\end{minipage} & \begin{minipage}[b]{\linewidth}\raggedright
runtime
\end{minipage} \\
\midrule\noalign{}
\endhead
\bottomrule\noalign{}
\endlastfoot
\texttt{probe\_ccm\_exact\_matrix.py} & the per-frequency exact matrix; zero-side validation; even/odd spectra at three windows; \(\epsilon_\lambda\) vs Fuchs; Rayleigh excess & \$\sim\$3 min \\
\texttt{probe\_quantified\_weld\_e2e.py} & min--max transfer on exact data; Mellin/strip transfer constants; the spectral realization tables above & \$\sim\$4 min \\
\texttt{probe\_fuchs\_second\_rung.py} & tower identification of the excited state; Courant--Fischer gap bound; Fuchs rung ratios from FD prolates & \$\sim\$4 min \\
\texttt{probe\_ccm\_seam\_weld.py} & Meixner--Schäfke rates; \(\widehat{k_\lambda}\to\Xi\) end-to-end; exact Poisson mirror and its finite-\(\lambda\) breakage & \$\sim\$3 min \\
\texttt{attack\_v2\_hardening.py} & stability audit: \(N\)-refinement of all spectra; quadrature doubling; 40 vs 60 digits; overlap deficit with random-subspace control; root-finder basin; finite-difference budget for \(r/\epsilon\) & \$\sim\$15 min \\
\end{longtable}
}

The same files, with their recorded outputs, are mirrored at: \textbf{{[}GitHub repository link{]}} --- provided so that nothing need be executed from an e-mail attachment.

\subsection{8. Closing remark}\label{closing-remark}

The verification reported here was carried out as part of a broader study of the quadratic form \(QW_\lambda\) --- a quantified, conditional formulation of the convergence strategy of {[}CCM{]}, Section 7, and a spectral analysis of the sector structure underlying the ladder of Section 4; in particular, a closed-form mechanism for the smallness of \(QW_\lambda(k_\lambda)\) (the two defining conditions of \(h_\lambda\) --- evenness and vanishing integral --- annihilate exactly the two pole terms of the Weil functional, placing the global \(\mathcal E\)-image in the radical of the global form and reducing the windowed Rayleigh quotient to a below-window tail) --- which we are preparing separately. We would be glad to share it, and we welcome any correction or comment on the present note. It is a pleasure to acknowledge the clarity of {[}CCM{]}: every formula needed for an independent implementation is stated there explicitly, which is what made this verification possible.

\subsection{References}\label{references}

\begin{enumerate}
\def\labelenumi{\arabic{enumi}.}
\tightlist
\item
  A. Connes, C. Consani, H. Moscovici, \emph{Zeta spectral triples}, arXiv:2511.22755 (2025). \textbf{{[}CCM{]}}
\item
  A. Connes, C. Consani, \emph{Spectral triples and \(\zeta\)-cycles}, Enseign. Math. \textbf{69} (2023), no. 1--2, 93--148; arXiv:2106.01715. \textbf{{[}CC-cycles{]}}
\item
  A. Connes, W. D. van Suijlekom, \emph{Quadratic forms, real zeros and echoes of the spectral action}, Commun. Math. Phys. \textbf{406} (2025), 312; arXiv:2511.23257. \textbf{{[}CS{]}}
\item
  A. Connes, \emph{The Riemann Hypothesis: past, present and a letter through time}, arXiv:2602.04022 (2026). \textbf{{[}L{]}}
\item
  A. Connes, C. Consani, \emph{Weil positivity and trace formula: the archimedean place}, Selecta Math. (N.S.) \textbf{27} (2021), no. 4, Paper No.~77; arXiv:2006.13771.
\item
  A. Connes, \emph{Trace formula in noncommutative geometry and the zeros of the Riemann zeta function}, Selecta Math. (N.S.) \textbf{5} (1999), no. 1, 29--106.
\item
  A. Connes, H. Moscovici, \emph{The UV prolate spectrum matches the zeros of zeta}, Proc. Natl. Acad. Sci. USA \textbf{119} (2022), no. 22.
\item
  A. Groskin, \emph{High-precision approximation of Riemann zeros via the truncated Weil form}, arXiv:2605.20224 (2026).
\item
  D. Śliwiński, \emph{Spectral analysis of the \(D^{(\lambda,N)}_{\log}\) operators}, arXiv:2601.12133 (2026).
\item
  W. H. J. Fuchs, \emph{On the eigenvalues of an integral equation arising in the theory of band-limited signals}, J. Math. Anal. Appl. \textbf{9} (1964), 317--330.
\item
  J. Meixner, F. W. Schäfke, \emph{Mathieusche Funktionen und Sphäroidfunktionen}, Springer (1954).
\item
  D. Slepian, H. O. Pollak, \emph{Prolate spheroidal wave functions, Fourier analysis and uncertainty I}, Bell Syst. Tech. J. \textbf{40} (1961), 43--63.
\end{enumerate}

\end{document}
